% ====================================================================
\chapter{Pruebas y validación}
\label{cap:pruebas_validacion}
% ====================================================================

En el presente capítulo se expone la estrategia de verificación y validación del sistema de traducción en tiempo real de Lengua de Señas Argentina (LSA). Se detallan el plan de pruebas adoptado, la ejecución de los experimentos sobre los distintos componentes y el análisis cuantitativo y cualitativo de los resultados obtenidos frente a los requerimientos del proyecto.

% ====================================================================
\section{Plan de pruebas}
\label{sec:plan_pruebas}
% ====================================================================

El plan de pruebas establece la estrategia sistemática para comprobar el correcto funcionamiento, el desempeño en tiempo real y la calidad del software desarrollado.

\subsection{Alcance y tipos de pruebas}
La evaluación del sistema abarca desde la recolección de datos y el entrenamiento de modelos de aprendizaje profundo hasta la integración completa de los componentes en una videollamada web. Se aplicaron los siguientes tipos de pruebas:
\begin{itemize}
    \item \textbf{Pruebas unitarias y de componentes:} Verificación del extractor de puntos clave (\textit{landmarks}), la lógica de filtrado gestual (\texttt{RepeatGate}), los clasificadores de señas y la cuantización de los modelos de lenguaje.
    \item \textbf{Pruebas de integración:} Evaluación de las interfaces entre la extensión de navegador (Manifest V3), el motor de inferencia local, las llamadas a API en Google Cloud Run y el cliente de videoconferencia (Google Meet).
    \item \textbf{Pruebas de rendimiento y latencia:} Medición de tiempos de respuesta por etapa del canal de procesamiento (\textit{pipeline}), consumo de memoria en cliente y servidor, y tasa de fotogramas por segundo.
    \item \textbf{Pruebas de sistema de extremo a extremo (\textit{End-to-End}):} Validación del circuito completo de traducción (LSA a español y español a LSA) en entornos de llamada reales.
    \item \textbf{Pruebas de robustez y casos límite:} Evaluación del clasificador frente a cambios de iluminación, distorsiones ópticas de cámara web, presencia de múltiples personas en el fondo y variabilidad en la velocidad de articulación gestual.
\end{itemize}

\subsection{Entornos de evaluación}
Las pruebas se ejecutaron en dos entornos representativos:
\begin{enumerate}
    \item \textbf{Estación de desarrollo (GPU):} Utilizada para la evaluación estática fuera de línea, el entrenamiento de modelos y las pruebas de velocidad máxima de inferencia.
    \item \textbf{Equipo objetivo de cliente (Laptop sin GPU):} Configuración de gama media estándar orientada a verificar el comportamiento en tiempo real, el consumo de CPU/RAM y la latencia perceptible por el usuario final.
    \item \textbf{Infraestructura en la nube:} Instancia de Google Cloud Run (configuración de 4 vCPU y 8\,GiB de memoria RAM) conectada mediante HTTPS para los servicios remotos de traducción semántica y reproducción de avatar.
\end{enumerate}

\subsection{Criterios de aceptación y trazabilidad}
Cada caso de prueba se vinculó de manera directa con los requerimientos funcionales y no funcionales del sistema. En la Tabla~\ref{tab:trazabilidad_pruebas} se sintetizan los criterios de aceptación principales definidos para la validación.

\begin{table}[H]
    \centering
    \small
    \caption{Matriz de trazabilidad entre requerimientos y criterios de aceptación para pruebas.}
    \label{tab:trazabilidad_pruebas}
    \begin{tabularx}{\textwidth}{l X X}
        \toprule
        \textbf{Módulo} & \textbf{Requerimiento asociado} & \textbf{Criterio de aceptación} \\
        \midrule
        Clasificador Gestual & Clasificación en tiempo real de vocabulario de 100 señas LSA. & Exactitud Top-1 en cámara real $> 75\%$ y Top-3 $> 90\%$. Latencia de inferencia $< 10$\,ms. \\
        Traductor Semántico & Conversión bidireccional entre glosas y español fluido. & Exactitud sintáctica/semántica $> 85\%$, BLEU $> 90$, BERTScore $> 98$. \\
        Motor Local & Filtrado de repeticiones y control de estado de gestos. & Supresión de falsos positivos y rebotes sin descartar repeticiones intencionales. \\
        Integración Meet & Captura e inyección de video/texto sin alterar la videollamada. & Sin interrupción de la pista de video principal; renderizado de subtítulos sin pérdidas. \\
        \bottomrule
    \end{tabularx}
\end{table}

% ====================================================================
\section{Resultados de la verificación}
\label{sec:resultados_verificacion}
% ====================================================================

Esta sección presenta la evidencia empírica obtenida de la ejecución del plan de pruebas sobre cada componente del sistema.

\subsection{Evaluación del clasificador gestual (TinySkeletonClassifier)}
\label{sec:res_clasificador}

El modelo de clasificación gestual se evaluó tanto en entornos estáticos fuera de línea como en pruebas en vivo frente a cámara web en la estación local.

\subsubsection{Evolución del corpus y vocabulario}
La incorporación de clases al modelo avanzó de manera paulatina conforme se consolidaron las muestras del dataset propio de LSA. En la Tabla~\ref{tab:evolucion_corpus} se detalla la progresión del vocabulario y las muestras recopiladas.

\begin{table}[H]
    \centering
    \small
    \caption{Progresión cronológica del conjunto de datos y vocabulario de LSA.}
    \label{tab:evolucion_corpus}
    \begin{tabular}{lcccc}
        \toprule
        \textbf{Nro} & \textbf{Clases} & \textbf{Muestras/Clase} & \textbf{Muestras Entrenamiento} & \textbf{Muestras Validación} \\
        \midrule
        1 & 91 & 50 & 3\,640 & 910 \\
        2 & 94 & 60 & 4\,512 & 1\,128 \\
        3 & 97 & 60 & 4\,656 & 1\,164 \\
        4 & 100 & 60 & 4\,800 & 1\,200 \\
        5 & 100 & 100 & 8\,000 & 2\,000 \\
        \bottomrule
    \end{tabular}
\end{table}

\subsubsection{Comparativa entre evaluación fuera de línea y cámara en vivo}
En la Tabla~\ref{tab:brecha_offline_camara} se expone la diferencia observada entre la exactitud de validación sobre el dataset estático y el rendimiento real obtenido frente a la cámara web.

\begin{table}[H]
    \centering
    \small
    \caption{Exactitud de validación fuera de línea frente al acierto en cámara. Las filas de 91 a 97 clases corresponden a una toma por seña. Las de 100 clases corresponden a diez tomas por seña.}
    \label{tab:brecha_offline_camara}
    \begin{tabular}{llccc}
        \toprule
        \textbf{Checkpoint} & \textbf{Fecha} & \textbf{Val. Top-1} & \textbf{Cámara Top-1} & \textbf{Cámara Top-3} \\
        \midrule
        16 fotogramas, 91 clases & 29--30/07 & 97{,}36\,\% & 61{,}5\,\% & 83{,}5\,\% \\
        32 fotogramas, 128/4/2 & 31/07 & 96{,}48\,\% & 71{,}4\,\% & 87{,}9\,\% \\
        32 fotogramas, 64/2/1 & 31/07 & 97{,}36\,\% & 67{,}0\,\% & 90{,}1\,\% \\
        Optuna v1, 32 fotogramas & 31/07 & 97{,}69\,\% & 66{,}0\,\% & 84{,}6\,\% \\
        Baseline 16f, luz ambiental, 94 clases & 18/08 & 98{,}40\,\% & 67{,}0\,\% & 89{,}4\,\% \\
        Baseline 16f, iluminación dirigida, 94 clases & 28/08 & 98{,}40\,\% & 86{,}2\,\% & 95{,}7\,\% \\
        Baseline 16f, 97 clases & 15--17/09 & 98{,}28\,\% & 77{,}3\,\% & 94{,}8\,\% \\
        100 clases, 10 tomas, luz adecuada & 28/10 & 99{,}00\,\% & 99{,}6\,\% & 100{,}0\,\% \\
        100 clases, 10 tomas, luz reducida & 28--29/10 & 99{,}00\,\% & 94{,}6\,\% & 98{,}7\,\% \\
        \bottomrule
    \end{tabular}
\end{table}

En las iteraciones de 94 clases, el control de la iluminación del entorno subió el acierto en cámara del 67{,}0\,\% al 86{,}2\,\% en Top-1 sobre la misma arquitectura, sin modificar los pesos del modelo. Esa comparación no fija encuadre ni persona. La evaluación del checkpoint de 100 clases, descrita más adelante, sí conserva ambos y aísla la luz.

\subsubsection{Resolución temporal: 16 frente a 32 fotogramas}
La Tabla~\ref{tab:frames_julio} sintetiza los experimentos realizados para determinar la longitud óptima de la ventana temporal de entrada.

\begin{table}[H]
    \centering
    \small
    \caption{Evaluación de la resolución temporal en julio de 2026 (91 clases, 50 videos por clase).}
    \label{tab:frames_julio}
    \begin{tabular}{lcccc}
        \toprule
        \textbf{Iteración} & \textbf{Arquitectura} & \textbf{Val. Top-1} & \textbf{Cámara Top-1} & \textbf{Cámara Top-3} \\
        \midrule
        A, 16 fotogramas & 128 dim / 4 cabezas / 2 capas & 97{,}36\,\% & 61{,}5\,\% & 83{,}5\,\% \\
        B, 32 fotogramas & 128 dim / 4 cabezas / 2 capas & 96{,}48\,\% & 71{,}4\,\% & 87{,}9\,\% \\
        C, 32 fotogramas & 64 dim / 2 cabezas / 1 capa & 97{,}36\,\% & 67{,}0\,\% & 90{,}1\,\% \\
        \bottomrule
    \end{tabular}
\end{table}

A pesar de las ganancias iniciales observadas con 32 fotogramas, la incorporación de aumentaciones avanzadas de datos en agosto confirmó que la ventana de 16 fotogramas ofreció mayor capacidad de generalización sobre señas estáticas y redujo la redundancia temporal.

\subsubsection{Impacto de la corrección de relación de aspecto} % TODO: Esta seccion si queres la volamos a la reconcha de la lora, pero es un dato interesante que no se ve en ningun lado. La idea es que la red aprende a reconocer la seña en un marco 16:9, y si le das un 4:3 con bandas negras, la red se confunde. Esto es un hallazgo de la etapa de pruebas, no del entrenamiento. Ahre completaba el texto con Claude
La estandarización del marco de captura desde proporciones $4:3$ hacia $16:9$ con bandas de relleno generó un impacto positivo directo en la clasificación, según se detalla en la Tabla~\ref{tab:aspect}.

\begin{table}[H]
    \centering
    \small
    \caption{Efecto de la corrección de relación de aspecto sobre 36 tomas de letras y dígitos.}
    \label{tab:aspect}
    \begin{tabular}{lc}
        \toprule
        \textbf{Configuración de entrada} & \textbf{Exactitud Top-1} \\
        \midrule
        Captura original en relación 4:3 & 66{,}7\,\% \\
        Corrección analítica en eje $y$ ($\times 1{,}333$) & 88{,}9\,\% \\
        Corrección simulada en ejes $y$ y $z$ & 75{,}0\,\% \\
        \bottomrule
    \end{tabular}
\end{table}

La transformación del eje $y$ aportó un incremento de 22,2 puntos porcentuales en exactitud. Corrigio confusiones sistemáticas como el caso de la letra \texttt{I} clasificada previamente como \texttt{ojo}.

\subsubsection{Evaluación en cámara del modelo de 100 clases}
\label{sec:eval_100_camara}

El prototipo integrado carga el modelo \texttt{model\_new\_2026\_10\_20}: 100 clases, 16 fotogramas, 128 dimensiones, 4 cabezas y 2 capas. El entrenamiento usó 100 clips por clase (8\,000 secuencias de entrenamiento y 2\,000 de validación) y alcanzó una exactitud de validación del 99{,}0\,\%, con pérdida $0{,}028$.

La prueba en cámara consistió en dos sesiones sobre ese mismo archivo. La mano, el encuadre y el extractor se mantuvieron fijos. Cambió la iluminación de la habitación. La etiqueta de luz adecuada o reducida la anota quien ejecuta el script; el programa no mide la iluminancia. Cada sesión recorrió las 100 señas con diez tomas consecutivas (1\,000 inferencias por condición), en modo automático y con mano derecha. No hubo tomas omitidas.

La cámara se solicita a $1280\times 720$ y entrega $640\times 480$. El script rellena bandas hasta $853\times 480$, el mismo marco 16:9 de los videos de entrenamiento. El extractor es MediaPipe Holistic en navegador (\texttt{@mediapipe/holistic@0.5.1675471629}): pose y dos manos, sin malla facial, vector de 225 coordenadas y resampleo a 16 fotogramas. Entre inferencias hay un segundo de espera.

La Tabla~\ref{tab:eval_100_luz} resume el resultado. El criterio de aceptación del plan (Top-1 mayor que 75\,\% y Top-3 mayor que 90\,\%) se cumple en las dos condiciones de iluminación.

\begin{table}[H]
    \centering
    \small
    \caption{Evaluación en cámara del modelo (\texttt{model\_new\_2026\_10\_20}). Diez tomas por seña, misma persona y mismo encuadre.}
    \label{tab:eval_100_luz}
    \begin{tabular}{lcc}
        \toprule
        \textbf{Métrica} & \textbf{Luz adecuada} & \textbf{Luz reducida} \\
        \midrule
        Top-1 & 99{,}6\,\% (996/1000) & 94{,}6\,\% (946/1000) \\
        Top-3 & 100{,}0\,\% (1000/1000) & 98{,}7\,\% (987/1000) \\
        Señas con 10 aciertos sobre 10 & 97 & 77 \\
        Errores que restan & 4 & 54 \\
        Confianza media en aciertos & $0{,}979$ & $0{,}967$ \\
        Confianza media en errores & $0{,}632$ & $0{,}771$ \\
        \bottomrule
    \end{tabular}
\end{table}

Con luz adecuada, 97 señas registran diez aciertos en diez tomas. Restan cuatro errores, y en los cuatro la seña pedida ocupa el segundo lugar: \texttt{A} se confunde dos veces con \texttt{6}; \texttt{ahora\_hoy} una vez con \texttt{tuyo}; \texttt{repetir} una vez con \texttt{esposo\_a}. El subconjunto de 37 letras, eñe y dígitos alcanza un 99{,}5\,\% de Top-1.

Con luz reducida el Top-1 desciende 5,0 puntos y 20 señas dejan de completar el diez sobre diez. El Top-3 permanece por encima del 98\,\%. Las caídas más marcadas, de señas que en luz adecuada estaban en diez sobre diez o en nueve sobre diez, son \texttt{T} (3/10, hacia \texttt{I}), \texttt{como} (3/10, hacia \texttt{numero}), \texttt{años} (5/10, hacia \texttt{G}), \texttt{hola} (6/10, hacia \texttt{chau}) y \texttt{ahora\_hoy} (6/10, con destinos distintos). En trece tomas la seña pedida no entra al Top-3. En luz reducida la confianza media de los errores sube a $0{,}771$: la red se equivoca con más certeza que en luz adecuada.

Esta sesión es la evidencia del modelo integrado. El 77{,}3\,\% del 17/09 describe un checkpoint anterior, con una toma por seña y otro extractor. El criterio de selección del producto es el acierto en cámara; con el checkpoint empaquetado en ILSA, ese acierto queda por encima del umbral de RNF-03 en las dos condiciones de iluminación.

\subsection{Resultados de los módulos de traducción semántica}
\label{sec:res_traduccion}

La evaluación de los modelos de lenguaje reducidos (SLM) cuantizados en formato GGUF (\texttt{Q4\_K\_M}) se ejecutó sobre la partición final del corpus (119 oraciones para el canal directo LSA $\rightarrow$ Español y 102 oraciones para el canal inverso Español $\rightarrow$ LSA).

\subsubsection{Evaluación del canal directo (LSA a Español)}
En la Tabla~\ref{tab:sem_directo} se presentan las métricas cualitativas y cuantitativas para el canal directo.

\begin{table}[H]
\centering
\footnotesize
\caption{Canal directo (LSA $\rightarrow$ Español), resultados finales (119 oraciones de evaluación). Índices en escala 0--100.}
\label{tab:sem_directo}
\begin{tabular}{lrrrrrrr}
\toprule
\textbf{Modelo} & \textbf{Exacto} & \textbf{BLEU} & \textbf{ROUGE-L} & \textbf{METEOR} & \textbf{BERTScore} & \textbf{COMET} & \textbf{ms} \\
\midrule
Qwen2.5 0,5B & 92,44 & 95,58 & 95,78 & 95,51 & 99,16 & 97,47 & 141,05 \\
Qwen2.5 1,5B & 94,96 & 96,32 & 98,02 & 97,35 & 99,20 & 97,36 & 256,30 \\
Qwen2.5 3B   & 95,80 & 97,88 & 98,39 & 97,84 & 99,20 & 97,56 & 514,38 \\
Llama 3.2 1B & 89,92 & 92,64 & 95,94 & 96,13 & 98,58 & 96,83 & 231,28 \\
SmolLM2 1,7B & 90,76 & 93,53 & 96,48 & 95,97 & 98,85 & 97,12 & 438,31 \\
\bottomrule
\end{tabular}
\end{table}

\subsubsection{Evaluación del canal inverso (Español a LSA)}
En la Tabla~\ref{tab:sem_inverso} se detallan las métricas correspondientes al módulo de estructuración en glosas para el canal inverso.

\begin{table}[H]
\centering
\footnotesize
\caption{Canal inverso (Español $\rightarrow$ LSA), resultados finales (102 oraciones de evaluación). Índices en escala 0--100.}
\label{tab:sem_inverso}
\begin{tabular}{lrrrrrrr}
\toprule
\textbf{Modelo} & \textbf{Exacto} & \textbf{BLEU} & \textbf{ROUGE-L} & \textbf{METEOR} & \textbf{BERTScore} & \textbf{COMET} & \textbf{ms} \\
\midrule
Qwen2.5 0,5B & 87,25 & 91,29 & 96,20 & 91,25 & 98,86 & 94,28 & 155,97 \\
Qwen2.5 1,5B & 91,18 & 94,14 & 97,17 & 92,75 & 99,22 & 95,98 & 290,82 \\
Qwen2.5 3B   & 89,22 & 92,51 & 96,85 & 92,71 & 99,21 & 95,45 & 537,58 \\
Llama 3.2 1B & 89,22 & 92,66 & 96,36 & 92,89 & 99,13 & 95,29 & 245,70 \\
SmolLM2 1,7B & 77,45 & 82,90 & 91,54 & 86,31 & 97,68 & 92,37 & 458,39 \\
\bottomrule
\end{tabular}
\end{table}

El modelo Qwen2.5 1,5B alcanzó el equilibrio óptimo en ambas direcciones, con un acierto exacto de 94,96\,\% en canal directo y 91,18\,\% en canal inverso. Mantuvo tiempos de inferencia en CPU inferiores a 300\,ms.

\subsection{Pruebas de la lógica de control del motor local}
La batería de pruebas diseñada para validar la lógica del filtro \texttt{RepeatGate} y los tiempos de corte de enunciado produjo los resultados consignados en la Tabla~\ref{tab:bateria_logica}.

\begin{table}[H]
    \centering
    \small
    \caption{Resultados de las pruebas funcionales sobre la lógica de control del motor local.}
    \label{tab:bateria_logica}
    \begin{tabularx}{\textwidth}{l l X}
        \toprule
        \textbf{Caso de prueba} & \textbf{Entrada ingresada} & \textbf{Resultado emitido} \\
        \midrule
        Dígitos repetidos & \texttt{1 1 1 2 2} & Aceptación total: \texttt{1 1 1 2 2} \\
        Letra repetida tres veces & \texttt{A A A B} & Filtrado a máximo de dos: \texttt{A A B} \\
        Rebote del clasificador & \texttt{HOLA HOLA} & Descarte de duplicado inmediato: \texttt{HOLA} \\
        Repetición no consecutiva & \texttt{HOLA VER HOLA} & Conservación de la secuencia completa \\
        Glosa repetida descartada a los 3\,s & Glosa repetida & El cierre se desplaza a los 7\,s \\
        Movimiento continuo hasta 5\,s & Sin glosa nueva & Cierre de enunciado a los 9\,s \\
        Deletreo dactilológico & \texttt{J U A N} & Conversión directa a «Juan.» sin consulta a nube \\
        Deletreo con ambigüedad & \texttt{T I N A} & Derivación al traductor por coincidencia con \texttt{OJO} \\
        \bottomrule
    \end{tabularx}
\end{table}

\subsection{Tratamiento de fallas de integración en Google Meet}
Durante la evaluación sobre la interfaz de usuario en llamadas activas se identificaron inconvenientes técnicos que fueron solucionados mediante intervenciones específicas de arquitectura, como se detalla en la Tabla~\ref{tab:fallas_integracion}.

\begin{table}[H]
    \centering
    \small
    \caption{Anomalías detectadas en la integración con Google Meet y soluciones aplicadas.}
    \label{tab:fallas_integracion}
    \begin{tabularx}{\textwidth}{@{}>{\raggedright\arraybackslash}p{0.42\textwidth} >{\raggedright\arraybackslash}X@{}}
        \toprule
        \textbf{Falla detectada} & \textbf{Intervención de ingeniería aplicada} \\
        \midrule
        Bloqueo del inicio de video en Meet por intercepción permanente de cámara. & Activación del reemplazo de la pista únicamente con el motor ILSA encendido. \\
        \addlinespace[2pt]
        Pérdida de enlace al recargar la extensión con la pestaña de la reunión abierta. & Recarga automática de la pestaña de Meet ante el reinicio del contexto de la extensión. \\
        \addlinespace[2pt]
        Detección de falsas manos provocada por movimiento del torso. & Exigencia de un número mínimo de cuadros consecutivos con mano detectada antes de abrir el gesto. \\
        \addlinespace[2pt]
        Interferencia en la captura por movimiento de personas en el fondo. & Delimitación del área de actividad basada exclusivamente en los puntos articulares de las manos. \\
        \addlinespace[2pt]
        Sobrecarga de CPU en la transferencia de imágenes hacia la ventana \textit{offscreen}. & Reducción de la resolución de los cuadros transferidos a $240$ píxeles de ancho con compresión ajustada. \\
        \bottomrule
    \end{tabularx}
\end{table}

\subsection{Pruebas de latencia y perfilado end-to-end en llamada real}
\label{sec:eval_telemetria_vivo}

Con fecha 29 de octubre de 2026 se ejecutó una sesión de pruebas end-to-end sobre la extensión integrada en videollamada web activa. Durante dicha sesión se capturaron eventos de telemetría en ambos canales de traducción (directo e inverso), registrando la latencia por etapa, el consumo de memoria del proceso cliente (RSS) y la tasa de fotogramas por segundo (FPS).

En la Tabla~\ref{tab:telemetria_resumen} se sintetizan los valores promedios, mínimos y máximos recolectados durante las ejecuciones de prueba.

\begin{table}[H]
    \centering
    \small
    \caption{Resumen de métricas de perfilado y latencia registradas durante la prueba del 29/10/2026.}
    \label{tab:telemetria_resumen}
    \begin{tabularx}{\textwidth}{l X X X}
        \toprule
        \textbf{Métrica de ejecución} & \textbf{Canal Directo (LSA $\rightarrow$ Esp)} & \textbf{Canal Inverso (Esp $\rightarrow$ LSA)} & \textbf{Unidad} \\
        \midrule
        Eventos registrados & 9 & 7 & oraciones \\
        Pausa de cierre (configurada) & $4031{,}1 \pm 23{,}5$ & $0{,}0$ & ms \\
        Tiempo de inferencia remota & $1444{,}4$ ($896{,}7$ -- $2372{,}9$) & $1368{,}2$ ($826{,}6$ -- $2202{,}8$) & ms \\
        Tiempo de transporte de red & $268{,}2$ ($273{,}4$ -- $336{,}4$) & $253{,}3$ ($254{,}2$ -- $328{,}2$) & ms \\
        Tiempo de entrega / renderizado & $0{,}3$ ($0{,}2$ -- $0{,}6$) & $869{,}3$ ($105{,}9$ -- $1387{,}6$) & ms \\
        Latencia total de procesamiento & $1712{,}9$ ($1233{,}3$ -- $2674{,}2$) & $2490{,}8$ ($105{,}9$ -- $3612{,}1$) & ms \\
        Consumo de memoria RAM (RSS) & $168{,}8$ ($96{,}0$ -- $224{,}9$) & $231{,}2$ ($219{,}8$ -- $235{,}2$) & MB \\
        Tasa de procesamiento (Holistic) & $5{,}8$ -- $12{,}1$ & N/A & FPS \\
        Tasa de renderizado cliente & $45{,}0$ -- $60{,}0$ & $58{,}2$ -- $60{,}5$ & FPS \\
        Cumplimiento RNF-01 ($< 2{,}5$\,s) & 77{,}8\,\% (7/9) & 85{,}7\,\% (6/7) & aciertos \\
        \bottomrule
    \end{tabularx}
\end{table}

El análisis cualitativo y cuantitativo de estos datos permite extraer las siguientes observaciones técnicas:

\begin{enumerate}
    \item \textbf{Latencia de procesamiento y transporte:} El tiempo medio de inferencia en el servidor remoto oscila alrededor de los $1{,}4$ segundos para ambos canales. El tiempo de transporte por red de la consulta cifrada desde el cliente web hacia Google Cloud Run promedia $268$ ms en el canal directo y $253$ ms en el canal inverso, evidenciando una comunicación ligera y estable.
    \item \textbf{Diferencial de entrega entre canales:} En el canal directo la entrega es casi instantánea ($0{,}3$ ms) debido a que la respuesta del LLM se renderiza como subtítulo de texto sobre la interfaz. En el canal inverso, la fase de entrega demanda un promedio de $869{,}3$ ms (con picos de $1387{,}6$ ms) motivado por la secuenciación, decodificación e interpolación de los tensores de articulación del avatar tridimensional.
    \item \textbf{Tasa de fotogramas y fluidez de captura:} La extracción de puntos clave mediante MediaPipe Holistic en el navegador cliente operó de forma continua entre $5$ y $12$ FPS en la ventana de captura gestual, suficiente para reconstruir tensores de 16 fotogramas sin perder componentes finas del gesto. La interfaz de usuario conservó un refresco fluido constante de $60$ FPS en el canal inverso.
\end{enumerate}

% ====================================================================
\section{Métricas de calidad}
\label{sec:metricas_calidad}
% ====================================================================

La evaluación objetiva del sistema se consolidó mediante métricas cuantificables agrupadas en cuatro dimensiones principales: desempeño en tiempo real, eficiencia en el uso de recursos, calidad semántica y confiabilidad.

\subsection{Desglose de latencias y tiempos de respuesta}
\label{sec:latencias_pipeline}
En la Tabla~\ref{tab:latencias_etapas} se desglosan los componentes del retardo total en el canal directo LSA a español.

\begin{table}[H]
    \centering
    \small
    \caption{Desglose de los componentes del retardo de procesamiento en el canal directo.}
    \label{tab:latencias_etapas}
    \begin{tabularx}{\textwidth}{@{}>{\raggedright\arraybackslash}p{0.32\textwidth} >{\raggedright\arraybackslash}p{0.20\textwidth} >{\raggedright\arraybackslash}X@{}}
        \toprule
        \textbf{Etapa del pipeline} & \textbf{Duración / Retardo} & \textbf{Naturaleza del parámetro} \\
        \midrule
        Detección de fin de seña & $\approx 0{,}4$\,s & Configurado (12 fotogramas sin manos). \\
        Inferencia del clasificador & $0{,}78$\,ms (CPU) / $1{,}08$\,ms (GPU) & Medición en CPU/GPU sobre modelo de 16f. \\
        Separación mínima entre señas & $1{,}0$\,s & Configurado (reposo inter-gestual). \\
        Pausa de cierre de enunciado & $4{,}0$\,s & Configurado (inactividad para disparar traducción). \\
        % TODO: Falta viaje al servidor (ida y vuelta)
        Generación de oración (LLM) & $256{,}30$\,ms & Medido en CPU local sobre Qwen2.5 1,5B. \\
        Reevaluación de prompt (CPU) & $\approx 8$\,s & Ocurre únicamente ante pérdida de caché $KV$. \\
        Arranque en frío del servicio & $1$--$3$\,min & Ocurre si la instancia escala a cero en Cloud Run. \\
        Permanencia de subtítulo & $8{,}0$\,s & Configurado (visibilidad en interfaz). \\
        \bottomrule
    \end{tabularx}
\end{table}


\subsection{Consumo de recursos y perfilado en cliente}
\label{sec:perfilado_cliente_recursos}

A partir de los registros telemetry recopilados en cliente, se analizó el impacto del motor de inferencia local sobre los recursos de cómputo de la laptop objetivo sin GPU dedicada.

En la Tabla~\ref{tab:perfilado_recursos_log} se comparan los requerimientos de memoria y CPU observados durante la ejecución continua de la extensión.

\begin{table}[H]
    \centering
    \small
    \caption{Métricas de uso de memoria física (RSS) y uso de CPU en la estación cliente durante la prueba activa.}
    \label{tab:perfilado_recursos_log}
    \begin{tabular}{lcccc}
        \toprule
        \textbf{Estado del canal} & \textbf{RAM Promedio (MB)} & \textbf{RAM Máxima (MB)} & \textbf{CPU Promedio (\%)} & \textbf{CPU Máximo (\%)} \\
        \midrule
        Canal Directo (LSA $\rightarrow$ Esp) & 168{,}81 & 224{,}90 & 2{,}94\,\% & 8{,}30\,\% \\
        Canal Inverso (Esp $\rightarrow$ LSA) & 231{,}19 & 235{,}20 & 33{,}80\,\% & 106{,}20\,\% \\
        \bottomrule
    \end{tabular}
\end{table}

El análisis de perfilado muestra que:
\begin{itemize}
    \item \textbf{Eficiencia en memoria:} El canal directo presentó una huella de memoria sumamente baja, reduciéndose progresivamente desde los 224{,}9 MB iniciales hasta un régimen permanente de 96{,}0 MB. El canal inverso mantuvo un consumo ligeramente superior pero acotado, estabilizado en $\approx 231$ MB debido al almacenamiento en caché de las animaciones del avatar.
    \item \textbf{Carga del procesador:} Durante el procesamiento del canal directo el uso de CPU cliente se mantuvo extremadamente bajo (2{,}94\,\% promedio), demostrando la liviandad del clasificador gestual de 1{,}66 MB. En el canal inverso, el renderizado gráfico del avatar generó un incremento medio del uso de CPU hasta el 33{,}80\,\%, con picos transitorios de 106{,}20\,\% durante la decodificación inicial del audio entrante por el motor Vosk, retornando de inmediato al estado de reposo.
\end{itemize}
El factor principal de espera radica en la pausa de 4,0 segundos, parametrizada para garantizar que la persona sorda haya finalizado la oración antes de enviar la secuencia al modelo de lenguaje.

%TODO: Comienza con «La verificación de RNF-01 se ejecutó sobre el motor ILSA» y termina con «no registran memoria RSS ni porcentaje de CPU.»
La verificación de RNF-01 se ejecutó sobre el motor ILSA en Python, no sobre el ejecutable empaquetado, en una sesión de Google Meet contra Cloud Run. El volcado reúne 11 emisiones del canal directo y 18 del inverso. Cinco de esas 29 no invocaron el SLM (entradas de diccionario o de una sola palabra) y no entran al cómputo del requerimiento. De las 24 que sí lo invocaron, 21 (87\,\%) respetaron el techo de 2{,}0\,s de inferencia y 14 (58\,\%) el objetivo interno de 1{,}5\,s. La meta de RNF-01 es el techo de 2\,s, no el objetivo de 1{,}5\,s.

\begin{table}[H]
    \centering
    \small
    \caption{Inferencia del SLM en Cloud Run. Solo emisiones que invocaron el modelo. El techo de RNF-01 es 2{,}0\,s.}
    \label{tab:rnf01_sesion}
    \begin{tabular}{lccccc}
        \toprule
        \textbf{Canal} & \textbf{$n$ con IA} & \textbf{Mediana} & \textbf{Media} & \textbf{$\leq 1{,}5$\,s} & \textbf{$\leq 2{,}0$\,s} \\
        \midrule
        Directo & 9 & 1{,}35\,s & 1{,}71\,s & 5 & 6 \\
        Inverso & 15 & 1{,}48\,s & 1{,}41\,s & 9 & 15 \\
        Total & 24 & --- & --- & 14 & 21 \\
        \bottomrule
    \end{tabular}
\end{table}

En el canal directo, los tres incumplimientos del techo coinciden con enunciados de 7, 4 y 10 glosas (2{,}03\,s, 2{,}54\,s y 3{,}46\,s). El máximo es un DNI de diez tokens. Una emisión adicional (DNI de cinco dígitos, 1{,}92\,s) quedó entre 1{,}5 y 2{,}0\,s. El canal inverso no superó 2{,}0\,s en ninguna de las 15 emisiones con modelo.

\begin{table}[H]
    \centering
    \small
    \caption{Cinco relojes medidos en la sesión integrada. RNF-01 cubre solo la columna de inferencia.}
    \label{tab:cinco_relojes_medidos}
    \begin{tabularx}{\textwidth}{@{}>{\raggedright\arraybackslash}p{0.22\textwidth} >{\raggedright\arraybackslash}p{0.28\textwidth} >{\raggedright\arraybackslash}X >{\raggedright\arraybackslash}p{0.12\textwidth}@{}}
        \toprule
        \textbf{Reloj} & \textbf{Canal directo} & \textbf{Canal inverso} & \textbf{RNF-01} \\
        \midrule
        Pausa & 4{,}00--4{,}14\,s en las 11 emisiones & 0 (la oración llega cerrada) & No \\
        Cola & 0\,ms & 0\,ms & No \\
        Inferencia & mediana 1{,}35\,s (9 con IA) & mediana 1{,}48\,s (15 con IA) & Sí \\
        Transporte & $\approx 0{,}30$\,s & $\approx 0{,}30$\,s; un pico de 0{,}75\,s & No \\
        Entrega & subtítulo $< 1$\,ms & avatar entre $\approx 0$ y 2{,}4\,s & No \\
        \bottomrule
    \end{tabularx}
\end{table}

La heurística de cierre de 4{,}0\,s se observó en las 11 emisiones directas y se informa aparte de RNF-01. La cola nula corresponde a un solo cliente con el servidor libre: no prueba concurrencia. El transporte no explica los tres fallos de inferencia. El subtítulo del canal directo es instantáneo; la espera perceptible en ese canal es la pausa de diseño. La entrega del avatar (descarga de archivos \texttt{.pose}) es ajena a RNF-01; con el archivo ya en disco, esa etapa cae a milisegundos.

En el canal directo de esa sesión, el clasificador local registró una mediana de 4{,}6\,ms sobre 48 señas. La extracción Holistic en el navegador registró una media de 7 fotogramas por segundo. El lienzo hacia Meet registró una mediana de 52 fps en Directo y de 60 fps en Inverso. Estos volcados no registran memoria RSS ni porcentaje de CPU.

\subsection{Consumo de memoria y huella del sistema}
La cuantización en formato GGUF de 4 bits permitió reducir la huella de memoria del modelo de lenguaje a valores compatibles con despliegues eficientes. En la Tabla~\ref{tab:pesos_cuant} y la Tabla~\ref{tab:ram_inferencia} se sintetizan el tamaño de los archivos y el consumo estimado de memoria RAM en inferencia con un contexto de 2048 tokens.

% TODO: Chequear estos valores
\begin{table}[H]
\centering
\footnotesize
\caption{Dimensiones de los archivos cuantizados y compresión alcanzada.}
\label{tab:pesos_cuant}
\begin{tabular}{lrrrr}
\toprule
\textbf{Modelo} & \textbf{16 bits (GiB)} & \textbf{GGUF (MiB)} & \textbf{GGUF (GiB)} & \textbf{bits/parámetro} \\
\midrule
Qwen2.5 0,5B & 0,913 & 379,4 & 0,370 & 6,5 \\
Qwen2.5 1,5B & 2,869 & 940,4 & 0,918 & 5,1 \\
Qwen2.5 3B   & 5,756 & 1\,840,5 & 1,797 & 5,0 \\
Llama 3.2 1B & 2,291 & 770,3 & 0,752 & 5,3 \\
SmolLM2 1,7B & 3,185 & 1\,006,7 & 0,983 & 4,9 \\
\bottomrule
\end{tabular}
\end{table}

\begin{table}[H]
\centering
\footnotesize
\caption{Estimación del consumo de memoria RAM en inferencia (contexto de 2048 tokens).}
\label{tab:ram_inferencia}
\begin{tabular}{lrrr}
\toprule
\textbf{Modelo} & \textbf{Caché \(KV\) (GiB)} & \textbf{16 bits + caché (GiB)} & \textbf{GGUF + caché (GiB)} \\
\midrule
Qwen2.5 0,5B & 0,023 & 0,936 & 0,394 \\
Qwen2.5 1,5B & 0,055 & 2,923 & 0,973 \\
Qwen2.5 3B   & 0,070 & 5,826 & 1,867 \\
Llama 3.2 1B & 0,062 & 2,354 & 0,815 \\
SmolLM2 1,7B & 0,375 & 3,560 & 1,358 \\
\bottomrule
\end{tabular}
\end{table}

En la réplica de producción desplegada en Cloud Run, ambos canales (directo e inverso) con el modelo Qwen2.5 1,5B consumen de manera conjunta $1{,}946$\,GiB de memoria RAM, situándose holgadamente dentro del límite asignado de $8$\,GiB.

\subsection{Comparativa de rendimiento según arquitectura de hardware}
En la Tabla~\ref{tab:hw_comparacion} se evalúa la viabilidad técnica del sistema según el hardware disponible para justificar la selección de una arquitectura híbrida cliente-servidor.

\begin{table}[H]
    \centering
    \small
    \caption{Comparación del comportamiento del sistema entre la estación de desarrollo y el equipo cliente.}
    \label{tab:hw_comparacion}
    \begin{tabularx}{\textwidth}{@{}>{\raggedright\arraybackslash}p{0.25\textwidth} >{\raggedright\arraybackslash}X >{\raggedright\arraybackslash}X@{}}
        \toprule
        \textbf{Componente} & \textbf{Estación de desarrollo (GPU)} & \textbf{Equipo objetivo (Laptop sin GPU)} \\
        \midrule
        Extracción MediaPipe & Procesamiento fluido en tiempo real. & Tasa de fotogramas reducida; requirió reajuste de umbrales a milisegundos. \\
        Clasificador local & Inferencia en $< 2$\,ms. & Inferencia en $< 5$\,ms sin impacto perceptible. \\
        Traductor local & Respuesta fluida. & Alta latencia en CPU, incompatible con la fluidez de la conversación. \\
        \bottomrule
    \end{tabularx}
\end{table}

La asignación de la clasificación al cliente de borde (modelo de $1{,}66$\,MB con 414\,309 parámetros) y la traducción al servidor en la nube garantizó el cumplimiento de los requerimientos de tiempo real en equipos sin aceleración por GPU.